\subsection{停机时间的赋值、Walsh 断点与账本平台的三重同一化}

正文小节 \ref{subsubsec:conclusion-godel-leyang-halting-time-dichotomies}
已经分别给出停机时间的 \(2\)-进代数化、Hausdorff 维数二值律、Fourier--Walsh 首断点、有限分辨率不可分辨下界、局部因子误差壁垒以及账本缺陷的断点剖面。
把这些接口并读后，可把同一停机时间参数重新压缩为一条完全同步的赋值--谱--账本链。

\begin{theorem}[停机时间的赋值--Walsh--账本三重同一化]\label{thm:derived-halting-valuation-walsh-ledger-unification}
\leanverified{paper\_derived\_halting\_valuation\_walsh\_ledger\_unification}
固定 \(B=2^L\)。对任意图灵机--输入对 \((M,x)\)，定义
\[
\alpha_{M,x}:=\sum_{n\ge 0}t_n(M,x)\,B^n\in\ZZ_2,
\]
其中 \(t_n(M,x)=1\) 当且仅当 \(M(x)\) 在不超过 \(n\) 步内停机。
令
\[
D_{M,x}:=\{0,\alpha_{M,x}e\}\subset T_2(E_{\min}),
\]
\[
k_\ast(M,x):=
\inf\Bigl\{
k\ge 0:\ \exists\,\chi,\ \mathrm{cond}(\chi)=k,\
\widehat{\mu}_{D_{M,x}}(\chi)\neq 1
\Bigr\},
\]
并记账本缺陷剖面为
\[
\kappa_m:=\mathbb E\bigl[\log|\pi_m^{-1}(Y)|\bigr].
\]
若 \(N\in\NN\cup\{\infty\}\) 为首次停机时间，则下列条件等价：
\[
N<\infty
\Longleftrightarrow
\alpha_{M,x}=-B^N
\Longleftrightarrow
v_2(\alpha_{M,x})=LN
\Longleftrightarrow
k_\ast(M,x)=LN+1
\Longleftrightarrow
\kappa_m=\min(m,N)\log 2\ \text{并最终平台化}.
\]
并且：
\begin{enumerate}
\normalfont
  \item 若 \(N<\infty\)，则
  \[
  \dim_H K_{D_{M,x}}=\dim(\mu_{D_{M,x}})=\frac{1}{L},
  \qquad
  \mu_{D_{M,x}}\ \text{非原子};
  \]
  \item 若 \(N=\infty\)，则
  \[
  K_{D_{M,x}}=\{0\},
  \qquad
  \mu_{D_{M,x}}=\delta_0,
  \qquad
  \kappa_m=m\log 2,
  \qquad
  k_\ast(M,x)=\infty.
  \]
\end{enumerate}
因此，停机时间在该 Gödel--Lee--Yang 地址化中同时显形为一个 \(2\)-进赋值、一个 Fourier--Walsh 首断点以及一条纤维账本曲线的平台起点。
\end{theorem}

\begin{proof}
由引理 \ref{lem:conclusion-halting-alpha-algebraic-form}，
\[
N=\infty \iff \alpha_{M,x}=0,
\qquad
N<\infty \iff \alpha_{M,x}=-B^N=-2^{LN},
\]
故前两重等价与赋值公式
\(
v_2(\alpha_{M,x})=LN
\)
立即成立。
第三重等价由定理 \ref{thm:conclusion-halting-walsh-breakpoint} 给出：若 \(N<\infty\)，则全部导子 \(k\le LN\) 上 Walsh 系数都等于 \(1\)，而存在导子 \(LN+1\) 首次偏离；若 \(N=\infty\)，则该首断点不存在。
第四重等价由命题 \ref{prop:conclusion-halting-ledger-defect-kink} 给出，其断点剖面正是
\[
\kappa_m=\min(m,N)\log 2.
\]
最后，几何与测度分岔由定理 \ref{thm:conclusion-halting-dimension-binary-law}
及推论 \ref{cor:conclusion-halting-measure-atomic-branch} 直接给出。证毕。
\end{proof}

\begin{proposition}[有限分辨率盲区与局部因子误差壁垒的同步化]\label{prop:derived-halting-resolution-localfactor-synchronization}
沿定理 \ref{thm:derived-halting-valuation-walsh-ledger-unification} 的记号。
对任意固定分辨率 \(m\ge 1\)，任何仅依赖截断观测
\[
X\bmod 2^m\ZZ_2 e
\]
的随机化学习器，都不可能在所有 \((M,x)\) 上以统一有界错误率区分“停机”与“不停机”；更精确地，只要某个停机实例满足
\[
N\ge \Bigl\lceil\frac{m}{L}\Bigr\rceil,
\]
则其截断分布与不停机情形完全一致。

另一方面，Lee--Yang 局部因子端的相应可计算构造 \(g_{M,x}\in SL_2(\CC)\) 满足
\[
g_{M,x}=I \iff N=\infty,
\]
并且对任意固定 \(s>0\) 与素数 \(p\)，若 \(N<\infty\) 则
\[
\Bigl|L_p(s,g_{M,x})-(1-p^{-s})^{-2}\Bigr|
\le C(s,p)\,2^{-N}.
\]
故同一停机时间 \(N\) 同时控制两类屏障：在截断学习侧，它给出一个严格的不可分辨盲区；在局部因子侧，它给出一个指数小的平凡模型误差壁垒。
\end{proposition}

\begin{proof}
第一条正是定理 \ref{thm:conclusion-halting-finite-resolution-indistinguishability}
的内容：当 \(N\ge \lceil m/L\rceil\) 时，停机实例在模 \(2^m\) 截断后与不停机情形同分布。
第二条由定理
\ref{thm:conclusion-halting-langlands-parameter-degeneration-barrier}
立即给出。二者共享同一时间标尺 \(N\)，只是分别表现为“完全不可区分”与“指数接近平凡”，故构成同步壁垒。证毕。
\end{proof}

\begin{remark}[赋值、频谱与账本的共同坐标]\label{rem:derived-halting-threefold-coordinate}
定理 \ref{thm:derived-halting-valuation-walsh-ledger-unification} 与命题
\ref{prop:derived-halting-resolution-localfactor-synchronization}
表明：停机时间在这里并不只是递归论意义上的 first-hit 参数，而是同一个隐藏坐标在 \(2\)-进赋值、Walsh 谱首断点、分形/测度相变、纤维账本平台与局部因子误差壁垒五个界面上的同步显影。
\end{remark}

\endinput
